\section{Holomorphic families, residue traces and the suspended eta invariant}
\label{tr:sec}

\providecommand{\trOp}{\operatorname{Op}}
\providecommand{\trOpc}{\operatorname{Op}_c}
\providecommand{\trang}[1]{\langle #1\rangle}
\providecommand{\trbint}{\overline{\int}}
\providecommand{\trres}{\operatorname{res}}
\providecommand{\trTrt}{\widetilde{\Tr}}
\providecommand{\trRN}[1]{\overline{\R^{#1}}}
\providecommand{\trP}{\mathcal P}
\providecommand{\trad}{\operatorname{ad}}

This appendix proves the background items (B3), (B7) and (B8) of
Section~\ref{subsec:background}. We follow the route of
\cite[Chs.~6 and 15]{MelroseMicro}. Instead of the complex powers
\(Q^{-s}\) we use \emph{any} holomorphic family \(Q(s)\) of order \(-s\) with
\(Q(0)=\mathrm{Id}\). Such families are written down explicitly, and nothing
about the spectral theory of \(Q\) is needed. The items are proved in the
following forms.
\begin{itemize}
\item[(B3\('\))] Proposition~\ref{tr:prop:Q} constructs an admissible family
\(Q(s)\) in the cusp calculus. Its indicial family is translation
invariant. Theorem~\ref{tr:thm:mero} proves the two-variable meromorphy of
\(\Tr(A(s,z)x^z)\) for holomorphic families \(A(s,z)\) of order \(m-s\),
with simple poles on the lines \(s=m+n-k\), \(z=1-p-j\), and identifies the
residues. Proposition~\ref{tr:prop:transfer} lists which results of
Sections~\ref{sec:residues}--\ref{sec:index} hold with \(Q(s)\) in place of
\(Q^{-s}\). It also explains why the final results do not depend on the
family.
\item[(B7\('\))] Theorem~\ref{tr:thm:bTr}: \(\bTr\) is a trace on
\(\Psis^{\Z}(H)\). Theorem~\ref{tr:thm:susreg} computes the
regularization of \(\bTr\) by a holomorphic family; this corrects
\cite[Prop.~8]{MN}. Theorem~\ref{tr:thm:Dirac}: for operators in APS
normal form,
\(\eta(A_0)=(\pi i)^{-1}\overline{\int}\Tr(A_0^{-1}\partial_\tau A_0)\,d\tau\)
is the APS eta invariant \(\eta_{B}(0)\). Proposition~\ref{tr:prop:etabar}
relates the functional \(\etab\) of \cite[(89)]{MN} to \(\eta(A_0)\).
\item[(B8\('\))] Theorem~\ref{tr:thm:res}: the residue is a trace. It is
given by the local formula and does not depend on the coordinates or on the
regularizing family. Proposition~\ref{tr:prop:unique} proves uniqueness for
\(n\ge2\) (for traces vanishing on smoothing operators).
\end{itemize}
The quoted inputs are these. From (B1) we use the local structure of the
cusp and suspended calculi: composition, symbols, the Taylor series at the
front face, and inverses of invertible elliptic suspended families. From
\cite{Notes} we use the small-time heat expansion of \(\Tr(Be^{-tB^2})\)
and the regularity of \(\eta_B(s)\) at \(s=0\). Only
Proposition~\ref{tr:prop:etabar}(ii) uses complex powers, and only on the
boundary. Every quoted step is marked.

Two statements need correction; we prove the corrected forms, which the
main text states.
\begin{enumerate}
\item The regular value at \(s=0\) of \(\bTr(BQ_0^{-s})\) is \emph{not} in
general \(\bTr(B)\), as claimed in \cite[Prop.~8]{MN}. There
is a local correction term (Theorem~\ref{tr:thm:susreg},
Example~\ref{tr:ex:counter}; corrected statement
Proposition~\ref{prop:Tr8}). It propagates to
Corollary~\ref{cor:Tr61}(2) and Proposition~\ref{prop:eta}; see
Correction~\ref{tr:corr:Prop8}.
\item If \(Q^{-s}\) is replaced by a general family \(Q(s)\), then
Proposition~\ref{prop:IF} acquires an extra coboundary term,
\(\frac12\Trds([A,B]\kappa_Q)\) (Remark~\ref{rem:IFkappa}). This term vanishes
when \([A,B]\in\cI\), so the index theorem is unaffected
(Proposition~\ref{tr:prop:transfer}).
\end{enumerate}

\subsection{Symbols of complex order, traces and Seeley's computation}
\label{tr:subsec:prelim}

Let \(\trRN N\) be the radial compactification of \(\R^N\). Let
\(\rho\in\CI(\trRN N)\) be a boundary defining function that is positive on
\(\R^N\) and equal to \(|\zeta|^{-1}\) for \(|\zeta|\ge1\).

\begin{definition}\label{tr:def:sym}
Let \(U\subset\R^k\) be open and \(\mu\in\C\). A \emph{classical symbol of
order \(\mu\)} on \(U\times\R^N\) is a function \(a=\rho^{-\mu}\beta\) with
\(\beta\in\CI(U\times\trRN N;M(r,\C))\). The Taylor series of \(\beta\) in
\(1/|\zeta|\) at the sphere at infinity gives the expansion
\(a\sim\sum_{j\ge0}a_{\mu-j}\), where
\(a_{\mu-j}(u,\zeta)=|\zeta|^{\mu-j}\beta_j(u,\zeta/|\zeta|)\) is homogeneous
of degree \(\mu-j\) and \(|a-\sum_{j<J}a_{\mu-j}|\le C_J|\zeta|^{\operatorname{Re}
\mu-J}\) for \(|\zeta|\ge1\), locally uniformly in \(u\). Let
\(\Omega\subset\C^l\) be open and \(\mu:\Omega\to\C\) holomorphic. A
\emph{holomorphic family of order \(\mu(w)\)} is
\(a(w)=\rho^{-\mu(w)}\beta(w)\) with \(\beta\in\CI(\Omega\times U\times
\trRN N)\) holomorphic in \(w\).
\end{definition}

A change of \(\rho\) multiplies \(\beta(w)\) by a holomorphic power of a
positive smooth function. So the definition does not depend on \(\rho\).

\begin{lemma}[Seeley's computation \cite{Seeley}]\label{tr:lem:polar}
Let \(a(w)\) be a holomorphic family of order \(\mu(w)\), with compact
support in \(u\). Then \(I(w;u)=\int_{\R^N}a(w;u,\zeta)\,d\zeta\) is
holomorphic on \(\{\operatorname{Re}\mu(w)<-N\}\). For every \(J\in\N\),
\[
 I(w;u)=\sum_{j<J}\frac{1}{j-N-\mu(w)}\int_{S^{N-1}}a_{\mu(w)-j}(w;u,\theta)
 \,d\theta+I_J(w;u),
\]
where \(I_J\) is holomorphic on \(\{\operatorname{Re}\mu(w)<J-N\}\) and smooth
in \(u\). In particular, if \(\mu(w)=m-s\) with \(w=(s,z')\), then \(I\) is
meromorphic in \(w\) with at most simple poles on the hyperplanes
\(s=m+N-j\). The residue there is \(\int_{S^{N-1}}a_{-N}\,d\theta\), where
\(a_{-N}\) is the term of homogeneity \(-N\) of the symbol \(a(w)\) at
\(s=m+N-j\).
\end{lemma}

\begin{proof}
The integral over \(|\zeta|\le1\) is holomorphic for all \(w\). For
\(|\zeta|\ge1\) write \(\zeta=\theta/r\), \(r\in(0,1]\), so that
\(d\zeta=r^{-N-1}dr\,d\theta\) and \(a=r^{-\mu}\beta(w;u,r,\theta)\). Taylor's
formula gives \(\beta=\sum_{j<J}r^j\beta_j+r^J\beta_{(J)}\), with all terms
smooth and holomorphic. For \(\operatorname{Re}\mu<-N\),
\(\int_0^1r^{j-\mu-N-1}dr=(j-N-\mu)^{-1}\). The remainder
\(\int_{S^{N-1}}\int_0^1r^{J-\mu-N-1}\beta_{(J)}\,dr\,d\theta\) is
holomorphic for \(\operatorname{Re}\mu<J-N\). Since
\(a_{\mu-j}(\theta)=\beta_j(\theta)\) on \(|\theta|=1\), this proves the
formula. The expressions for different \(J\) agree where both are defined,
so they define the continuation.
\end{proof}

\begin{lemma}\label{tr:lem:tc}
Let \(M\) be a closed manifold of dimension \(N\) and \(E\to M\) a bundle.
\begin{enumerate}
\item If \(A\in\Psi^\mu(M;E)\) with \(\operatorname{Re}\mu<-N\), then \(A\)
is of trace class, its kernel is continuous, and
\(\Tr A=\int_M\tr K_A|_\Delta\). If \(A=\phi\trOp(a)\psi\) in a chart, with
\(\phi,\psi\in\CI_c(U)\), then
\(\Tr A=(2\pi)^{-N}\int\!\!\int\phi\psi\tr a\,d\zeta\,du\).
\item If \(A\in\Psi^{\mu_1}\), \(B\in\Psi^{\mu_2}\) and
\(\operatorname{Re}(\mu_1+\mu_2)<-N\), then \(\Tr(AB)=\Tr(BA)\).
\end{enumerate}
\end{lemma}

\begin{proof}
(Adapted from \cite[Ch.~15]{MelroseMicro}.) (1) With a partition of unity,
\(A\) is a finite sum of operators \(\phi\trOp(a)\psi\) plus a smoothing
operator. For a smoothing operator the statement is standard. Write
\(a=a'\trang\zeta^{\mu/2}\), with \(a'\) of order \(\mu/2\). In the left
quantization \(\trOp(a')\trang D^{\mu/2}=\trOp(a)\) exactly. The operator
\(\trang D^{\mu/2}\psi\) is Hilbert--Schmidt, because \(\trang
\zeta^{\operatorname{Re}\mu/2}\in L^2\). The operator \(\phi\trOp(a')\) is
Hilbert--Schmidt, because its kernel has \(L^2\) norm in \(u'\) at most
\(\|a'(u,\cdot)\|_{L^2}\). So \(A\) is of trace class. Its trace is the
pairing of the two kernels. Since \(a\) is integrable in \(\zeta\), this
pairing is \(\int K_A(u,u)\,du\).

(2) Regard the operators as elements of the H\"ormander classes
\(\Psi^{\operatorname{Re}\mu}_{\mathrm H}\). The argument of (1) bounds the
trace norm of \(\phi\trOp(a)\psi\) by \(C\sup|a|\trang\zeta^{-\nu}\) for
\(\nu<-N\). So \(\Tr\) is continuous on \(\Psi^\nu_{\mathrm H}\), \(\nu<-N\).
Cut off the local symbols of \(A\) to \(\chi(\epsilon\zeta)a\), keeping the
part of \(A\) away from the diagonal. This gives smoothing operators
\(A_\epsilon\to A\) in \(\Psi^{\operatorname{Re}\mu_1+\delta}_{\mathrm H}\)
for every \(\delta>0\). Choose \(\delta\) with
\(\operatorname{Re}(\mu_1+\mu_2)+\delta<-N\). Then \(\Tr(AB)=\lim\Tr(A_\epsilon
B)=\lim\Tr(BA_\epsilon)=\Tr(BA)\).
\end{proof}

\subsection{The Wodzicki residue: item (B8)}\label{tr:subsec:res}

Let \(M\) be closed, of dimension \(N\). (In (B8), \(N=n\).)

\begin{definition}\label{tr:def:holM}
Let \(\Omega\subset\C^l\) be open and \(\mu\) holomorphic on \(\Omega\). A
family \(A(w)\in\Psi^{\mu(w)}(M;E)\) is \emph{holomorphic of order \(\mu\)}
if the following hold.
\begin{enumerate}
\item For \(\chi\in\CI(M^2)\) vanishing near the diagonal, \(\chi K_{A(w)}\)
is holomorphic with values in \(\CI(M^2)\).
\item In every chart over which \(E\) is trivial, and for \(\phi,\psi\in
\CI_c(U)\), one has \(\phi A(w)\psi=\phi\trOp(a(w))\psi+R(w)\). Here \(a(w)\)
is a holomorphic family of classical symbols of order \(\mu(w)\), and
\(R(w)\) is a holomorphic family of smoothing operators.
\end{enumerate}
An \emph{admissible family} is an entire family \(Q(s)\) of order \(-s\)
with \(Q(0)=\mathrm{Id}\) whose principal symbol is scalar, \(q_s\,
\mathrm{Id}_E\).
\end{definition}

\begin{lemma}\label{tr:lem:holcalc}
\begin{enumerate}
\item Products of holomorphic families are holomorphic families, and the
orders add.
\item (Division.) If \(F(s)\) is holomorphic of order \(\mu(s)\) and
\(F(s_0)=0\), then \(F(s)=(s-s_0)F_1(s)\), with \(F_1\) holomorphic of order
\(\mu(s)\). In two variables: if \(F(s,z)\) vanishes on \(\{s=0\}\cup
\{z=0\}\), then \(F=szF_2\). If \(F(0,0)=0\), then \(F=sF_3+zF_4\). All
the \(F_i\) are holomorphic of order \(\mu\).
\item If the principal symbol of a holomorphic family of order \(\mu\)
vanishes identically, it is a holomorphic family of order \(\mu-1\).
\end{enumerate}
\end{lemma}

\begin{proof}
(1) is the composition theorem with holomorphic parameters. The composition
integral and the Borel summation of the symbol expansion can be carried out
locally uniformly in \(w\) \cite[Ch.~15]{MelroseMicro}. We do not repeat
the argument.

(2) Away from the diagonal, divide the holomorphic kernel. Near the
diagonal, \(\phi F(s_0)\psi=0\) forces \(\rho^{-\mu(s_0)}\beta(s_0)\in
S^{-\infty}\), so \(\beta(s_0)\) vanishes to infinite order at the sphere at
infinity. Then \(S(s)=\phi\trOp(\rho^{-\mu(s)}\beta(s_0))\psi\) is a
holomorphic \emph{smoothing} family. Replace \(\beta(s)\) by
\(\beta(s)-\beta(s_0)\) and \(R(s)\) by \(R(s)+S(s)\). Both now vanish at
\(s_0\), and both are divisible by \(s-s_0\) by Cauchy's formula. For two
variables: if \(F(0,z)=0\), then \(F=sF_1\) with \(z\) a parameter. Then
\(F_1(s,0)=0\) for \(s\ne0\), hence for all \(s\), and \(F_1=zF_2\). If
\(F(0,0)=0\), apply the one-variable statement to \(F(s,z)-F(0,z)\) and to
\(F(0,z)\).

(3) We have \(\beta|_{\rho=0}=0\), so \(\beta=\rho\beta'\).
\end{proof}

\begin{proposition}\label{tr:prop:QM}
Choose charts \(U_i\) trivializing \(E\), a subordinate partition of unity
\(\phi_i\), and \(\psi_i\in\CI_c(U_i)\) with \(\psi_i=1\) near
\(\supp\phi_i\). Then
\[
 Q_M(s)=\sum_i\phi_i\trOp\bigl(\trang\zeta^{-s}\bigr)\psi_i ,
 \qquad\trang\zeta=(1+|\zeta|^2)^{1/2},
\]
is an admissible family. Its principal symbol is
\(q_s=\sum_i\phi_i|\zeta|_{(i)}^{-s}\), where \(|\zeta|_{(i)}\) is the
Euclidean length in the coordinates of \(U_i\).
\end{proposition}

\begin{proof}
Each term is an entire family of order \(-s\). Since \(\trOp(1)=
\mathrm{Id}\), we get \(Q_M(0)=\sum_i\phi_i\psi_i=\sum_i\phi_i=\mathrm{Id}\).
\end{proof}

Complex powers \(Q^{-s}\) of a positive elliptic \(Q\in\Psi^1\), when they
are available, are also admissible. We do not need them.

\begin{theorem}[The residue trace]\label{tr:thm:res}
Let \(A\in\Psi^m(M;E)\), \(m\in\Z\). Let \(A(s)\) be a holomorphic family of
order \(m-s\) on a connected open set \(\Omega\ni0\) that meets
\(\{\operatorname{Re}s>m+N\}\), with \(A(0)=A\). For instance
\(A(s)=AQ(s)\) with \(Q\) admissible. Then \(\Tr A(s)\) extends from
\(\{\operatorname{Re}s>m+N\}\cap\Omega\) to a meromorphic function on
\(\Omega\). It has at most simple poles, at \(s\in m+N-\N_0\), and
\begin{equation}\label{tr:eq:res}
 \Res_{s=0}\Tr A(s)=\TrR(A):=(2\pi)^{-N}\sum_i\int_{U_i}\phi_i(u)
 \int_{|\zeta|=1}\tr a^{(i)}_{-N}(u,\zeta)\,dS(\zeta)\,du ,
\end{equation}
where \(a^{(i)}_{-N}\) is the term of homogeneity \(-N\) of the full symbol
of \(A\) in the chart \(U_i\). Consequently:
\begin{enumerate}
\item \(\TrR(A)\) does not depend on the family \(A(s)\), in particular not
on \(Q\).
\item (Coordinate invariance.) The density
\(\bigl(\int_{|\zeta|=1}\tr a_{-N}(u,\zeta)\,dS\bigr)|du|\) on \(M\) does not
depend on the chart or on the trivialization of \(E\). On \(S^*M\) it is
the fibre integral of \(\tr a_{-N}\,|\iota_R(\omega^N/N!)|\), where \(R\)
is the radial vector field. This is the formula (B8) and \eqref{eq:wodzicki}.
\item (Trace.) \(\TrR([A,B])=0\) for all \(A,B\in\Psi^{\Z}(M;E)\).
\item \(\TrR\) vanishes on \(\Psi^{m}\) for \(m<-N\), in particular on
\(\Psi^{-\infty}\).
\item Let \(Q\) be admissible. Then \([A,Q(s)]=sM_A(s)\), where \(M_A\) is
holomorphic of order \(m-1-s\). The map
\(D_QA:=M_A(0)=\frac{d}{ds}[A,Q(s)]|_{s=0}\) is a derivation of order
\(-1\), and \(\TrR(D_QA)=0\).
\item Let \(\widehat\Tr_Q(A)\) be the regular value of \(\Tr(AQ(s))\) at
\(s=0\). Then \(\widehat\Tr_Q([A,B])=\TrR(A\,D_QB)\). If \(Q'=Q+sF\) is
another admissible family, then \(\widehat\Tr_{Q'}(A)=\widehat\Tr_Q(A)+
\TrR(AF(0))\).
\end{enumerate}
\end{theorem}

\begin{proof}
Write \(A(s)=\sum_i\phi_iA(s)\psi_i+\sum_i\phi_iA(s)(1-\psi_i)\). The second
sum is a holomorphic smoothing family whose kernel vanishes near the
diagonal. So its trace vanishes. By Definition~\ref{tr:def:holM} and
Lemma~\ref{tr:lem:tc}(1),
\(\Tr(\phi_iA(s)\psi_i)=(2\pi)^{-N}\int\phi_i(u)\int a^{(i)}(s;u,\zeta)\,
d\zeta\,du+\Tr R_i(s)\). Here \(\Tr R_i(s)\) is entire, and we have used
that \(\psi_i=1\) near \(\supp\phi_i\). Lemma~\ref{tr:lem:polar} gives the
continuation and the poles. At \(s=0\) the residue is the integral of
\(a^{(i)}_{-N}(0)\), the term of homogeneity \(-N\) of the symbol of
\(A(0)=A\). This proves \eqref{tr:eq:res} and (1).

(2) The left side of \eqref{tr:eq:res} is intrinsic. Apply it to \(\phi A\)
with \(\phi\in\CI_c(U_i\cap U_j)\) arbitrary. The two local densities then
have the same integral against every \(\phi\), so they agree. In
coordinates, \(\omega^N/N!=\pm d\zeta\wedge du\). Contracting with
\(R=\zeta\cdot\partial_\zeta\) and restricting to \(|\zeta|=1\) gives
\(\pm dS\,du\). For \(h\) homogeneous of degree \(-N\), the form
\(\iota_R(h\,d\zeta)\) is closed on \(\R^N\setminus0\). So the fibre integral
does not depend on the choice of section of \(T^*M\setminus0\to S^*M\).

(3) Let \(A\in\Psi^{m_1}\), \(B\in\Psi^{m_2}\) and let \(Q\) be admissible. For
\(\operatorname{Re}s\) large, Lemma~\ref{tr:lem:tc}(2) gives
\(\Tr([A,B]Q(s))=\Tr(ABQ(s))-\Tr(AQ(s)B)=\Tr(A[B,Q(s)])\). The family
\([B,Q(s)]\) vanishes at \(s=0\), so it equals \(sM_B(s)\) by
Lemma~\ref{tr:lem:holcalc}. Hence \(\Tr([A,B]Q(s))=s\Tr(AM_B(s))\) has no
pole at \(0\).

(4) There is no pole if \(m<-N\).

(5) The principal symbol of \([A,Q(s)]\) of order \(m-s\) is
\([\sigma_m(A),q_s\mathrm{Id}]=0\). By Lemma~\ref{tr:lem:holcalc}(3),(2),
\([A,Q(s)]=sM_A(s)\) with \(M_A\) of order \(m-1-s\). Differentiating
\([AB,Q]=[A,Q]B+A[B,Q]\) shows that \(D_Q\) is a derivation. Finally,
\(\TrR(D_QA)=\Res_{s=0}\Tr(M_A(s))\) by \eqref{tr:eq:res}. For
\(\operatorname{Re}s\) large, \(\Tr(M_A(s))=s^{-1}\Tr([A,Q(s)])=0\).

(6) By the proof of (3), \(\Tr([A,B]Q(s))=s\Tr(AM_B(s))\). Its value at
\(0\) is \(\Res_{s=0}\Tr(AM_B(s))=\TrR(AM_B(0))\). For the last statement,
\(\Tr(AQ'(s))=\Tr(AQ(s))+s\Tr(AF(s))\), and \(F(0)\) has order \(0\).
\end{proof}

Formally \(Q(s)=e^{-s\lQ}\), so \(\lQ=-Q'(0)\) and \(D_QA=[A,Q'(0)]=
[\lQ,A]\). This is the derivation of \cite[\S6]{MN} and of
Lemma~\ref{lem:closeddefect}. Part (6) is that lemma. The principal symbol
of \(D_QA\) is \(-i\{\log q,a\}\) with
\(\log q:=-\partial_sq_s|_{s=0}\), at least when the full symbol of \(Q\)
is scalar in the chosen trivializations. In general a matrix commutator
term of lower order may be present \cite[Ch.~15]{MelroseMicro}.

\begin{remark}\label{tr:rem:sign}
Melrose regularizes with families \(E(s)\) of order \(+s\) and obtains
\(-\TrR\) \cite[Ch.~15]{MelroseMicro}. With the convention of order \(-s\)
used here and in \cite{MN}, the residue \(\TrR\) is \emph{plus} the
Wodzicki--Guillemin residue \cite{Wodzicki-residue,Guillemin-residue}.
\end{remark}

\begin{proposition}[Uniqueness]\label{tr:prop:unique}
Let \(M\) be connected and \(N\ge2\). Every trace \(\tau\) on
\(\Psi^{\Z}(M)\) that vanishes on \(\Psi^{-\infty}(M)\) is a multiple of
\(\TrR\).
\end{proposition}

\begin{proof}
Work in a coordinate ball \(U\). For a symbol \(a\) with \(u\)-support in a
compact set \(K\subset U\), put \(P_a=\trOp(a)\psi\), with \(\psi\in
\CI_c(U)\) equal to \(1\) near \(K\). Let \(\tilde\psi\in\CI_c(U)\) be equal
to \(1\) near \(\supp\psi\). In the left quantization, \([u_k,\trOp(a)]=
i\trOp(\partial_{\zeta_k}a)\) and \([D_{u_k},\trOp(a)]=\trOp(D_{u_k}a)\)
exactly. Hence
\begin{equation}\label{tr:eq:comm}
 [\tilde\psi u_k,P_a]=iP_{\partial_{\zeta_k}a},\qquad
 [\tilde\psi D_{u_k},P_a]\equiv P_{D_{u_k}a}\pmod{\Psi^{-\infty}} .
\end{equation}
The error in the second relation is \(\trOp(a)(D_{u_k}\psi)\). Its kernel
vanishes near the diagonal, so it is smoothing.

\emph{Step 1: degrees \(j\ne-N\).} Let \(p_j\) be homogeneous of degree
\(j\in\Z\) with \(u\)-support in \(K\), and let \(\chi(\zeta)\) be equal to
\(0\) near \(0\) and to \(1\) for \(|\zeta|\ge1\). By Euler's identity,
\(\sum_k\partial_{\zeta_k}(\zeta_k\chi p_j)=(j+N)\chi p_j+r\), where \(r\)
has compact support in \(\zeta\). By \eqref{tr:eq:comm}, \(P_{\chi p_j}\) is
a sum of commutators modulo \(\Psi^{-\infty}\), so \(\tau(P_{\chi p_j})=0\).
The same applies to symbols \(q\in S^{-N-1}\). Write \(q\sim\sum_{j\le-N-1}
q_j\) and let \(b_k\) be an asymptotic sum of the symbols
\(-i(j+N)^{-1}\zeta_k\chi q_j\). Then \(P_q-\sum_k[\tilde\psi u_k,P_{b_k}]\in
\Psi^{-\infty}\), so \(\tau(P_q)=0\).

\emph{Step 2: degree \(-N\).} Let \(p\) be homogeneous of degree \(-N\) with
\(u\)-support in \(K\). Put \(\varrho(u)=\int_{S^{N-1}}p(u,\theta)\,d\theta\).
Since \(N\ge2\), the sphere is connected, and we can solve
\(\Delta_{S}w(u,\cdot)=p(u,\cdot)-\varrho(u)/|S^{N-1}|\) on \(S^{N-1}\),
smoothly in \(u\) and with compact support in \(u\). Extend \(\nabla_Sw\) to
a vector field \(g\) that is tangent to the spheres and homogeneous of degree
\(1-N\). Then \(\operatorname{div}_\zeta g=|\zeta|^{-N}\Delta_Sw\). So
\(p-\varrho(u)|\zeta|^{-N}/|S^{N-1}|=\sum_k\partial_{\zeta_k}g_k\). Next fix
\(\chi_0\in\CI_c(U)\) with \(\int\chi_0=1\) and put \(\nu=\int_U\varrho\,
du\). The function \(\varrho-\nu\chi_0\) has compact support and integral
\(0\) on the ball \(U\), so it equals \(\sum_k\partial_{u_k}b_k\) with
\(b_k\in\CI_c(U)\). By \eqref{tr:eq:comm},
\[
 \tau(P_{\chi p})=\nu\,\tau(P_{\chi g_0}),\qquad
 g_0=\chi_0(u)|\zeta|^{-N}/|S^{N-1}| .
\]
Since \(\TrR(P_{\chi p})=(2\pi)^{-N}\nu\), this gives
\(\tau(P_{\chi p})=c_U\TrR(P_{\chi p})\) with
\(c_U=(2\pi)^N\tau(P_{\chi g_0})\).

\emph{Step 3: globalization.} With a partition of unity and Steps 1--2,
\(\tau(A)=\sum_ic_{U_i}\TrR(\phi_iA\psi_i)\). If \(U_i\cap U_j\neq
\emptyset\), choose an operator localized in a small ball inside
\(U_i\cap U_j\) with non-zero residue. This shows \(c_{U_i}=c_{U_j}\). Since
\(M\) is connected, all the constants agree.
\end{proof}

For \(N=1\) the sphere \(S^0\) has two points. A homogeneous function of
degree \(0\) on \(\R\setminus0\) has derivative \(0\), and there are two
independent residues, one for each sign of \(\xi\). This is
Correction~\ref{corr:n1}. Wodzicki shows that every trace vanishes on
\(\Psi^{-\infty}\) (quoted: \cite{Wodzicki-residue}). We do not use this.

\subsection{Holomorphic families in the cusp calculus: item (B3)}
\label{tr:subsec:cusp}

Fix a collar \([0,\epsilon)_x\times\partial X\) and put \(t=1/x\). We use
(B1) in the following form (cf.\ Sections~\ref{subsec:double}
and~\ref{subsec:symbols}). Let \(V\subset\R^{n-1}\) be a chart of
\(\partial X\) over which \(E\) is trivial, and let \(\phi,\psi\) be
supported in \([0,\epsilon)\times V\), with \(\psi=1\) near \(\supp\phi\). If
\(A\in x^p\Psic^m(X;E)\), then
\begin{equation}\label{tr:eq:Opc}
\begin{gathered}
 \phi A\psi=\phi\,\trOpc(x^pb)\,\psi+R,\\
 \trOpc(a)u(t,y)=(2\pi)^{-n}\!\int\! e^{i(t-t')\tau+i(y-y')\cdot\eta}
 a(1/t,y,\tau,\eta)\,u(t',y')\,dt'\,dy'\,d\tau\,d\eta .
\end{gathered}
\end{equation}
Here \(b\) is a classical symbol of order \(m\) in \(\zeta=(\tau,\eta)\), with
coefficients smooth in \((x,y)\in[0,\epsilon)\times V\), and
\(R\in x^p\Psic^{-\infty}(X;E)\). Interior charts have the usual
description, and \(A-\sum_i\phi_iA\psi_i\in x^p\Psic^{-\infty}\). The kernel
of an element of \(x^p\Psic^{-\infty}\) restricted to the diagonal is \(x^p g\,
x^{-2}|dx\,dy|\) with \(g\) smooth up to \(x=0\). The Taylor series
\eqref{eq:taylor} of \(A\) at the front face is computed from the Taylor
series of \(x^pb\) in \(x\) (powers of \(x\) on the left). The coefficient
\(A_{(j)}\in\Psis^m(\partial X)\) has local symbol equal to the coefficient
of \(x^j\) in \(x^pb\). Finally, if \(C\in x^p\Psic^\mu\) with
\(\operatorname{Re}\mu<-n\) and \(\operatorname{Re}z>1-p\), then \(Cx^z\) is of
trace class and \(\Tr(Cx^z)=\int_Xx^z\tr K_C|_\Delta\) (the proof of
Lemma~\ref{tr:lem:tc}(1) applies).

\begin{definition}\label{tr:def:cuspfam}
A family \(A(w)\in x^p\Psic^{\mu(w)}(X;E)\), \(w\in\Omega\subset\C^l\), is
\emph{holomorphic} if the following hold. In boundary charts
\eqref{tr:eq:Opc} holds with \(b=b(w)\) a holomorphic family of classical
symbols of order \(\mu(w)\), with coefficients smooth up to \(x=0\), and
with \(R=R(w)\) holomorphic in the Fr\'echet space \(x^p\Psic^{-\infty}\). The
analogous statement holds in interior charts. And \(A(w)-\sum_i\phi_iA(w)
\psi_i\) is holomorphic in \(x^p\Psic^{-\infty}\). An \emph{admissible
family} is an entire family \(Q(s)\in\Psic^{-s}(X;E)\) with
\(Q(0)=\mathrm{Id}\) and scalar principal symbol.
\end{definition}

\begin{lemma}\label{tr:lem:cuspcalc}
\begin{enumerate}
\item Products of holomorphic families are holomorphic. Orders and weights
add.
\item The division statements and part (3) of Lemma~\ref{tr:lem:holcalc}
hold verbatim.
\item If \(A(w)\) is holomorphic in \(x^p\Psic^{\mu(w)}\), then
\(x^zA(w)x^{-z}\) is holomorphic in \((w,z)\) in \(x^p\Psic^{\mu(w)}\), and
\(\partial_z(x^zAx^{-z})|_{z=0}=[\lx,A]\in x^{p+1}\Psic^{\mu-1}\).
\item Let \(Q_0(w)\in\Psis^{\mu(w)}(\partial X)\) be holomorphic
(Definition~\ref{tr:def:sus}). Let \(\omega,\tilde\omega\in
\CI_c([0,\epsilon))\) be equal to \(1\) near \(0\), with \(\tilde\omega=1\)
on \(\supp\omega\). Then \(\tilde\omega Q_0(w)\omega\) (acting in \((t,y)\))
is holomorphic in \(\Psic^{\mu(w)}\), and its Taylor series at the front
face is \(Q_0(w)\).
\item All the families act holomorphically on \(\dCI(X;E)\).
\end{enumerate}
\end{lemma}

\begin{proof}[Proof (sketch; (B1) with parameters)]
For (1), (4) and (5), the single-operator statements of (B1) hold with
holomorphic dependence on parameters, as in Lemma~\ref{tr:lem:holcalc}(1).
For (4): \(\tilde\omega(x')-1\) and \(\omega(x)-1\) vanish to infinite order
at the front face, where \(x'/x\to1\). For (2), the proof of
Lemma~\ref{tr:lem:holcalc}(2),(3) uses only the form of the symbols. For (3),
the kernel of \(x^zAx^{-z}\) is \((x/x')^zK_A\). Near the front face,
\((x/x')^z=(t'/t)^z=(1-xT)^z\) with \(T=t-t'\). This is smooth and entire
in \(z\) for \(|xT|<1\), and of polynomial growth in \(T\) elsewhere. The
kernel is rapidly decreasing in \(T\) there. The \(z\)-derivative at \(0\)
is \(\log(x/x')K_A\), which is \([\lx,A]\). Its leading term is computed in
Lemma~\ref{lem:twist}(3).
\end{proof}

\begin{proposition}[An explicit admissible family]\label{tr:prop:Q}
Let \(\{\chi_i(y)\}\) be a partition of unity on \(\partial X\) subordinate
to charts \(V_i\) over which \(E\) is trivial, and let \(\tilde\chi_i=1\)
near \(\supp\chi_i\). Let \(\omega,\tilde\omega\) be as in
Lemma~\ref{tr:lem:cuspcalc}(4). Let \(\{\phi_k\}\) be interior cut-offs
with \(\sum_i\omega\chi_i+\sum_k\phi_k=1\), and let \(\psi_k=1\) near
\(\supp\phi_k\). Put
\[
 Q(s)=\sum_i\omega\chi_i\,\trOpc\bigl(\trang\zeta^{-s}\bigr)\,
 \tilde\omega\tilde\chi_i+\sum_k\phi_k\trOp\bigl(\trang\xi^{-s}\bigr)\psi_k .
\]
Then \(Q\) is an admissible family in \(\Psic^{-s}(X;E)\). It is
translation invariant to all orders at \(\partial X\): its Taylor series at
the front face is
\[
 Q_0(s)=\sum_i\chi_i\,\trOp_{t,y}\bigl(\trang{(\tau,\eta)}^{-s}\bigr)\,
 \tilde\chi_i\in\Psis^{-s}(\partial X;E),\qquad Q_0(0)=\mathrm{Id}.
\]
\end{proposition}

\begin{proof}
The symbol \(\trang\zeta^{-s}=\rho^{s}\cdot(\rho\trang\zeta)^{-s}\) is an
entire family of order \(-s\) that does not depend on \((x,y)\). So each term
is an entire family by \eqref{tr:eq:Opc} and
Lemma~\ref{tr:lem:cuspcalc}(4). Its principal symbol is the positive scalar
\(\omega\sum_i\chi_i|\zeta|_{(i)}^{-s}+\sum_k\phi_k|\xi|_{(k)}^{-s}\).
Since \(\trOpc(1)=\mathrm{Id}\), we get \(Q(0)=\sum_i\omega\chi_i
\tilde\omega\tilde\chi_i+\sum_k\phi_k\psi_k=\mathrm{Id}\). The interior terms
have kernels supported away from the front face. The Taylor series of
the boundary terms is given by Lemma~\ref{tr:lem:cuspcalc}(4).
\end{proof}

\begin{proposition}[The derivation \(D_Q\)]\label{tr:prop:DQ}
Let \(Q\) be admissible and \(A\in x^p\Psic^m(X;E)\).
\begin{enumerate}
\item \([A,Q(s)]=sM_A(s)\), where \(M_A\) is holomorphic of order
\(m-1-s\) and weight \(x^p\). Put
\[
 D_QA:=M_A(0)=\frac{d}{ds}[A,Q(s)]\Big|_{s=0}=[\lQ,A],\qquad
 \lQ:=-Q'(0).
\]
Here \(Q'(0)\) is an operator on \(\dCI(X;E)\). It is not in the calculus.
\item \(D_Q\) is a derivation of \(x^{-\Z}\Psic^{\Z}(X;E)\). It maps
\(x^p\Psic^m\) into \(x^p\Psic^{m-1}\) and preserves \(\cI\), \(\Is\) and
\(\Id\). For genuine complex powers it is the derivation
\(\frac{d}{ds}Q^sAQ^{-s}|_{s=0}\) of Section~\ref{sec:derivations}.
\item If \(Q\) has Taylor series \(Q_0(s)\) at the front face, then
\(\In(D_QA)=D_{Q_0}\In(A)\), where \(D_{Q_0}B=\frac{d}{ds}[B,Q_0(s)]|_{s=0}\)
on \(\Psis^{\Z}(\partial X)\).
\item If \(\tilde Q\) is another admissible family, then \(\tilde Q-Q=sF(s)\)
with \(F(0)\in\Psic^0\). Hence
\[
 D_{\tilde Q}=D_Q+[L,\cdot\,],\qquad L=-F(0)=:\log\tilde Q-\lQ\in
 \Psic^0(X;E),
\]
and the two derivations differ by an inner derivation.
\item \(Q(s)Q(u)-Q(s+u)=su\,W(s,u)\), where \(W\) is holomorphic of order
\(-(s+u)\). The operator
\[
 \kappa_Q:=-W(0,0)=Q''(0)-Q'(0)^2\in\Psic^0(X;E)
\]
is in the calculus. For genuine complex powers \(\kappa_Q=0\).
\end{enumerate}
\end{proposition}

\begin{proof}
(1) and (2) are proved as Theorem~\ref{tr:thm:res}(5), using
Lemma~\ref{tr:lem:cuspcalc}(2). The ideals are preserved because \(M_A(s)\)
is a family of the same type as \(A\). (3) Take Taylor series at the front
face. (4) Use the division lemma. (5) Both \(Q(s)Q(u)\) and \(Q(s+u)\) are
holomorphic families of order \(-(s+u)\), and their difference vanishes on
\(\{s=0\}\cup\{u=0\}\). Apply \(\partial_s\partial_u\) at \((0,0)\) as
operators on \(\dCI\) (Lemma~\ref{tr:lem:cuspcalc}(5)).
\end{proof}

\begin{theorem}[Two-variable meromorphy]\label{tr:thm:mero}
Let \(m,p\in\Z\). Let \(\Omega\subset\C^2\) be a connected open set
containing points of \(\{\operatorname{Re}s>m+n,\ \operatorname{Re}z>1-p\}\),
and let \(A(s,z)\in x^p\Psic^{m-s}(X;E)\) be holomorphic on \(\Omega\). Then
\(F(s,z)=\Tr(A(s,z)x^z)=\Tr(x^zA(s,z))\) extends meromorphically to
\(\Omega\). Near every point it has the form
\(\sum_{\epsilon_1,\epsilon_2\in\{0,1\}}h_{\epsilon_1\epsilon_2}(s,z)\,
(s-s_k)^{-\epsilon_1}(z-z_j)^{-\epsilon_2}\) with \(h\) holomorphic,
\(s_k=m+n-k\) and \(z_j=1-p-j\), \(k,j\in\N_0\). Thus \(F\) has at most simple
poles on these lines, and \(sz\,F\) is holomorphic near \((0,0)\).
Moreover:
\begin{enumerate}
\item (Residue in \(s\).) \(\Res_{s=0}F(s,z)=(2\pi)^{-n}\int_{\Y}x^z\tr
a_{-n}(0,z)\,{}^c\nu\), where \(a_{-n}(0,z)\) is the term of homogeneity
\(-n\) of the cusp symbol of \(A(0,z)\), and \({}^c\nu\) is the density of
Lemma~\ref{lem:Trs} (\(x^{-2}dx\,dy\,dS\) in boundary charts). It depends
only on \(A(0,\cdot)\). For \(A(0,z)\in\Is\) it is \(\Trs(A(0,z)x^z)\).
\item (Residue in \(z\).) \(\Res_{z=0}F(s,z)\) depends only on \(A(s,0)\).
For \(\operatorname{Re}s>m+n\) it equals
\[
 \Res_{z=0}F(s,z)=\frac1{2\pi}\int_{\R}\Tr_{\partial X}\bigl(A(s,0)_{(1)}(\tau)
 \bigr)\,d\tau .
\]
\item (Double residue.) The coefficient of \((sz)^{-1}\) is
\(\Res_{z=0}\Res_{s=0}F=\Res_{s=0}\Res_{z=0}F=\Trds(A(0,0))\), the
functional of Lemma~\ref{lem:Trdsformula}.
\end{enumerate}
\end{theorem}

\begin{proof}
For \(\operatorname{Re}s>m+n\) and \(\operatorname{Re}z>1-p\),
\(F=\int_Xx^z\tr K_{A}|_\Delta=\sum\int x^z\tr K_{\phi_iA\psi_i}|_\Delta\),
because \(A-\sum\phi_iA\psi_i\) has kernel vanishing near the diagonal. The
interior terms are entire in \(z\) and are treated in \(s\) as in
Theorem~\ref{tr:thm:res}. In a boundary chart, by \eqref{tr:eq:Opc} and
\(dt=x^{-2}dx\),
\begin{gather*}
 \int x^z\tr K_{\phi A\psi}|_\Delta=\int_0^\epsilon\!\!\int_Vx^{z+p-2}
 D(s,z;x,y)\,dy\,dx,\\
 D=(2\pi)^{-n}\!\int\phi\tr b(s,z;x,y,\zeta)\,d\zeta+g(s,z;x,y),
\end{gather*}
with \(g\) smooth up to \(x=0\) and holomorphic. By
Lemma~\ref{tr:lem:polar}, \(D=\sum_{k<K}d_k(s,z;x,y)/(s-s_k)+D^{(K)}\), with
all terms smooth in \((x,y)\) up to \(x=0\), compactly supported, and
holomorphic in \((s,z)\). The term \(D^{(K)}\) is holomorphic for
\(\operatorname{Re}s>m+n-K\). For a smooth compactly supported \(h(x)\)
depending holomorphically on parameters,
\[
 \int_0^\epsilon x^{z+p-2}h(x)\,dx=\sum_{j<J}\frac{h^{(j)}(0)}{j!\,(z-z_j)}
 +(\text{holomorphic for }\operatorname{Re}z>z_J).
\]
Applying this to each term gives the local form of \(F\).

(1) The residue at \(s=0\) is the contribution of \(k=m+n\). It is the
integral over \((x,y)\) of
\[
 x^{z+p-2}(2\pi)^{-n}\int_{S^{n-1}}\phi\tr b_{-n}(0,z)\,dS .
\]
The term \(b_{-n}(0,z)\) is the homogeneous term of degree \(-n\) of the
symbol of \(A(0,z)\) (Lemma~\ref{tr:lem:polar}).

(2) The pole at \(z=0\) comes from \(j=1-p\). Its residue is the Taylor
coefficient of \(x^{1-p}\) in \(D(s,0)\), that is, the coefficient of
\(x^1\) in the diagonal density \(x^pD\) of \(A(s,0)\). For
\(\operatorname{Re}s>m+n\) the coefficient is
\((2\pi)^{-n}\int(x^pb)_{(1)}\,d\zeta=(2\pi)^{-1}\int[(2\pi)^{1-n}\int
(x^pb)_{(1)}\,d\eta]\,d\tau\), plus the corresponding coefficient of the
residual part. By (B1) this is the value on the diagonal of the convolution
kernel of \(A(s,0)_{(1)}(\tau)\), integrated over \(\tau\). Integration over
\(y\) gives the formula.

(3) In the local form, the coefficient of \((sz)^{-1}\) does not depend on
the order of the residues. Taking the \(z\)-residue of (1) picks the
coefficient of \(x^1\) in \(a_{-n}(0,0)\). This is the formula of
Lemma~\ref{lem:Trdsformula}.
\end{proof}

\begin{proposition}[Transfer to the functionals of MN]\label{tr:prop:transfer}
Let \(Q\) be admissible. Define \(Z(A;s,z)=\frac12\bigl(\Tr(Ax^zQ(s))+
\Tr(AQ(s)x^z)\bigr)\), the functionals \(\Trds\), \(\hTrd\), \(\hTrs\) by
\eqref{eq:44}, \(\widetilde Z(A;z)=Z(A;z,z)\), \(\Bf\) and
\(\IF=\IF_Q\) as in Sections~\ref{sec:boundaryindex}--\ref{sec:index}, and
\([\lQ,\cdot\,]:=D_Q\). Then:
\begin{enumerate}
\item Lemmas~\ref{lem:LR}, \ref{lem:Trds}, \ref{lem:Trs},
\ref{lem:Trdsformula}, \ref{lem:hTrs}, \ref{lem:Trd}, \ref{lem:hTrdQ},
\ref{lem:TrdsD}, \ref{lem:17}, \ref{lem:fedosov},
Corollary~\ref{cor:Tr61}(1), Proposition~\ref{prop:Bf}, and the closed-case
Lemma~\ref{lem:closeddefect} and Proposition~\ref{prop:closedindex} hold as
stated. In Lemma~\ref{lem:hTrdQ}, \(\log(Q'/Q)\) is replaced by
\(\log Q'-\lQ=-F(0)\), where \(Q'=Q+sF\).
\item Proposition~\ref{prop:IF} holds in the form
\[
 \IF_Q(A,B)=\tfrac12\bigl(\hTrd+\hTrs\bigr)\bigl(B\widetilde D(A)+
 \widetilde D(A)B\bigr)+\tfrac12\Trds\bigl([A,B]\kappa_Q\bigr).
\]
The last term is the Hochschild coboundary of the \(0\)-cochain
\(C\mapsto\frac12\Trds(C\kappa_Q)\) on \(\cA\). It vanishes if \([A,B]\in
\Is+\Id\), in particular if \([B]=[A]^{-1}\) in \(\cA\). It vanishes
identically for genuine complex powers. Hence Theorem~\ref{thm:index} and
its corollaries hold verbatim.
\item (Independence.) \(\Trds\) and \(\hTrs\) do not depend on \(Q\).
\(\hTrd\) changes by \(-\Trds(\,\cdot\,L)\) with \(L=\log Q'-\lQ\),
\(D_Q\) changes by the inner derivation \([L,\cdot\,]\), and \(\IF_Q\)
changes by a coboundary. Hence all Hochschild classes and the index are
independent of the family.
\end{enumerate}
\end{proposition}

The notes use the inverse powers \(Q^{s}\) at five places: Lemmas
\ref{lem:LR}, \ref{lem:Trds}, \ref{lem:TrdsD}, Proposition~\ref{prop:Bf}
and Proposition~\ref{prop:IF}. The proof replaces each conjugation by
\(Q^{\pm s}\) with a division identity. We need two lemmas.

\begin{lemma}\label{tr:lem:diag}
Let \(E(z)\in x^q\Psic^{k-z}\) be a holomorphic family near \(0\), and put
\(E^\flat=\frac{d}{dz}\bigl(E(z)-E(0)Q(z)\bigr)|_{z=0}\in x^q\Psic^k\). Then
\[
 z^2\Tr(E(z)x^z)=\Trds(E(0))+z\Bigl(\bigl(\hTrd+\hTrs\bigr)(E(0))+
 \Trds(E^\flat)\Bigr)+O(z^2).
\]
\end{lemma}

\begin{proof}
Write \(E(z)=E(0)Q(z)+zE_1(z)\), so that \(E_1(0)=E^\flat\). Then
\(\Tr(E(z)x^z)=\Tr(E(0)Q(z)x^z)+z\Tr(E_1(z)x^z)\). The first term is
\(Z^R(E(0);z,z)\). By Lemma~\ref{lem:LR} it differs from \(Z(E(0);z,z)\) by
\(\frac12z^2\Tr(E(0)P_1(z,z)x^z)\), and \(z^2\) times this is \(O(z^2)\).
Then use \eqref{eq:44} with \(s=z\). For the second term,
\(z^2\Tr(E_1(z)x^z)=\Trds(E_1(0))+O(z)\) by
Theorem~\ref{tr:thm:mero}(3), applied to the two-variable family
\((s,z)\mapsto E_1(s)\) and restricted to the diagonal.
\end{proof}

\begin{lemma}\label{tr:lem:commexp}
Let \(\mathcal Q(z)\in\Psic^{-z}\) be holomorphic with
\(\mathcal Q(0)=\mathrm{Id}\). Put \(L(z)=z^{-1}(x^zAx^{-z}-A)\) and
\(N(z)=z^{-1}[A,\mathcal Q(z)]\). Then
\[
 \Tr\bigl([A,B]\mathcal Q(z)x^z\bigr)=z\Tr\bigl(E(z)x^z\bigr),\qquad
 E=B\bigl(\mathcal QL-N\bigr),\qquad E(0)=B\widetilde D(A),
\]
where \(\widetilde D=[\lx,\cdot\,]-D_{\mathcal Q}\) and
\(D_{\mathcal Q}A=N(0)\).
\end{lemma}

\begin{proof}
For \(\operatorname{Re}z\) large, \(\Tr([A,B]\mathcal Qx^z)=
\Tr(B[\mathcal Qx^z,A])\) by the trace property. Moreover
\([\mathcal Qx^z,A]=\mathcal Q[x^z,A]+[\mathcal Q,A]x^z=z(\mathcal QL-N)x^z\).
\end{proof}

\begin{proof}[Proof of Proposition~\ref{tr:prop:transfer}]
(1) \emph{Lemma~\ref{lem:LR}.} The difference is \(Z^R-Z^L=\Tr(A[Q(s),x^z])\),
and \([Q(s),x^z]=P(s,z)x^z\) with \(P=Q(s)-x^zQ(s)x^{-z}\). The family \(P\)
vanishes for \(z=0\), and also for \(s=0\) because \(Q(0)=\mathrm{Id}\). So
\(P=sz\,P_1\) by Lemma~\ref{tr:lem:cuspcalc}(2), and
\(Z^R-Z^L=sz\Tr(AP_1x^z)\). The rest of the proof is verbatim.

\emph{Lemma~\ref{lem:Trds}.} We have
\(Z^L([A,B])=\Tr(A[B,x^zQ(s)])\) and \([B,x^zQ(s)]=x^zK(s,z)\) with
\(K=(x^{-z}Bx^z)Q(s)-Q(s)B\). Then \(K(0,0)=0\), so \(K=sK_1+zK_2\), and
\(sz\,Z^L([A,B])=s\cdot sz\Tr(C_1x^z)+z\cdot sz\Tr(C_2x^z)\) with
\(C_i=Ax^zK_ix^{-z}\). Both terms vanish at \((0,0)\) by
Theorem~\ref{tr:thm:mero}. For a change of family use \(Q'=Q+sF\) in place
of \(Q'^{-s}=(1+sC(s))Q^{-s}\). Vanishing on \(\Is+\Id\) follows from
Theorem~\ref{tr:thm:mero}(1),(2).

\emph{Lemmas~\ref{lem:Trs}, \ref{lem:Trdsformula}, \ref{lem:hTrs},
\ref{lem:Trd} and Corollary~\ref{cor:Tr61}(1).} These are
Theorem~\ref{tr:thm:mero}(1)--(3) together with the locality of the
residues. Independence of \(Q\) holds because a residue at \(s=0\) depends
only on the value of the family at \(s=0\). Theorem~\ref{tr:thm:res}
identifies \(\Trs\) with Wodzicki's residue on the double \(2X\).

\emph{Lemma~\ref{lem:hTrdQ}.} We have \(sz\,Z(A;Q')=sz\,Z(A;Q)+s\cdot
sz\,Z(AF(s))\). Comparing the coefficients of \(s\) gives
\(\hTrd(A;Q')=\hTrd(A;Q)+\Trds(AF(0))\).

\emph{Lemma~\ref{lem:TrdsD}.} Write \([B,Q(s)]=sM_B(s)\). Then
\(M_B(0)=[B,Q'(0)]=[\lQ,B]=D_QB\). By Theorem~\ref{tr:thm:mero}(3),
\(\Trds(D_QB)\) is the double residue of \(\Tr(M_B(s)x^z)\). For
\(\operatorname{Re}s\) large, \(s\Tr(M_B(s)x^z)=\Tr(BQ(s)x^z)-\Tr(Bx^zQ(s))
=\Tr(B[Q(s),x^z])=sz\Tr(BP_1(s,z)x^z)\). So \(sz\Tr(M_Bx^z)=z\cdot
sz\Tr(BP_1x^z)\), which vanishes at \((0,0)\). For
\([\lx,B]\), write \(x^wBx^{-w}=B+wL_B(w)\). Then
\(w\Tr(L_B(w)x^zQ(s))=\Tr(Bx^z(x^{-w}Q(s)x^w-Q(s)))\), and
\(x^{-w}Q(s)x^w-Q(s)=wsP'(w,s)\), because it vanishes at \(w=0\) and at
\(s=0\). So \(sz\Tr(L_B(0)x^zQ(s))=s\cdot sz\Tr(Bx^zP'(0,s))\), and the
double residue of \(L_B(0)=[\lx,B]\) vanishes.

\emph{Lemma~\ref{lem:17}.} In the first line,
\(\Trd([A,B]Q(s))=\Trd(A[B,Q(s)])=s\Trd(AM_B(s))\), with \(M_B\) as in the
previous paragraph. By Theorem~\ref{tr:thm:mero}(3) its value at \(s=0\) is
\(\Trds(AM_B(0))=\Trds(A[\lQ,B])\). The second line does not involve \(Q\).

\emph{Proposition~\ref{prop:Bf} and Lemma~\ref{lem:fedosov}.} Note that
\(\Tr([A,B]x^zQ(z))=\Tr([A,B]\mathcal Q_L(z)x^z)\) with
\(\mathcal Q_L=x^zQ(z)x^{-z}\), and that \(D_{\mathcal Q_L}=D_Q\) because
\(\mathcal Q_L'(0)=Q'(0)\). Apply Lemma~\ref{tr:lem:commexp} with
\(\mathcal Q=\mathcal Q_L\) and with \(\mathcal Q=Q\), and then
Lemma~\ref{tr:lem:diag}. This gives
\begin{align*}
 2\widetilde Z([A,B];z)&=\sum_{\mathcal Q\in\{\mathcal Q_L,Q\}}z\Tr
 \bigl(E_{\mathcal Q}(z)x^z\bigr)\\
 &=\frac{2\Trds(B\widetilde DA)}{z}+
 \sum_{\mathcal Q}\Bigl((\hTrd+\hTrs)(B\widetilde DA)+
 \Trds(E^\flat_{\mathcal Q})\Bigr)+O(z).
\end{align*}
So \(\Bf(A,B)=\Trds(B\widetilde DA)\), as in Proposition~\ref{prop:Bf}.
Lemma~\ref{lem:fedosov} only uses that \(\widetilde Z(C;\cdot)\) is entire
with value \(\Tr C\) at \(0\) for \(C\in\cI\). This holds because
\(Q(0)=\mathrm{Id}\).

(2) By the display above,
\(\IF_Q(A,B)=(\hTrd+\hTrs)(B\widetilde DA)+\frac12\Trds(E^\flat_{\mathcal
Q_L}+E^\flat_Q)\). We compute \(E^\flat=E'(0)+E(0)\ell\), with
\(\ell=\lQ=-Q'(0)\), as an identity of operators on \(\dCI\). Write
\(\trad_x=[\lx,\cdot\,]\), \(\trad_\ell=[\ell,\cdot\,]\) and
\(X=\trad_xA\). Then \(L(0)=X\), \(L'(0)=\frac12\trad_x^2A\),
\(N(0)=\trad_\ell A\) and \(N'(0)=\frac12[A,\mathcal Q''(0)]\). By
Proposition~\ref{tr:prop:DQ}(5), \(Q''(0)=\ell^2+\kappa_Q\) and
\(\mathcal Q_L''(0)=\ell^2+\kappa_Q-2[\lx,\ell]\). A direct computation
gives
\[
 E^\flat_Q=B\Bigl(\tfrac12\trad_x^2A-\trad_\ell\trad_xA+\tfrac12
 \trad_\ell^2A-\tfrac12[A,\kappa_Q]\Bigr),\qquad
 E^\flat_{\mathcal Q_L}=E^\flat_Q+B\bigl[A,[\lx,\ell]\bigr].
\]
Here we used \(-\frac12[A,\ell^2]-(\trad_\ell A)\ell=\frac12\trad_\ell^2A\).
By the Jacobi identity, \(\trad_x\trad_\ell-\trad_\ell\trad_x=
\trad_{[\lx,\ell]}\), and
\(\widetilde D=\trad_x-\trad_\ell\). Hence
\[
 E^\flat_Q+E^\flat_{\mathcal Q_L}=B\bigl(\widetilde D^2A-[A,\kappa_Q]\bigr).
\]
By Lemma~\ref{lem:TrdsD} and the trace property,
\(\Trds(B\widetilde D^2A)=-\Trds(\widetilde DB\,\widetilde DA)\). By
Lemma~\ref{lem:17},
\(\frac12(\hTrd+\hTrs)([\widetilde DA,B])=-\frac12\Trds(\widetilde DA\,
\widetilde DB)\). Combining,
\[
 \IF_Q(A,B)=\tfrac12(\hTrd+\hTrs)(B\widetilde DA+\widetilde DA\,B)
 -\tfrac12\Trds(B[A,\kappa_Q]),
\]
and \(-\Trds(B[A,\kappa_Q])=\Trds([A,B]\kappa_Q)\). The vanishing statements
hold because \(\Trds\) vanishes on \(\Is+\Id\supset\cI\).

(3) This follows from (1), (2), Proposition~\ref{tr:prop:DQ}(4) and
Lemma~\ref{lem:iD}.
\end{proof}

\subsection{The regularized trace on the suspended algebra: item (B7)}
\label{tr:subsec:sus}

Let \(H\) be a closed manifold of dimension \(d\). In the applications
\(H=\partial X\) and \(d=n-1\).

\begin{definition}\label{tr:def:sus}
For \(\mu\in\C\), \(\Psis^\mu(H;E)\) is the space of families
\(\R\ni\tau\mapsto C(\tau)\in\Psi^\mu(H;E)\) of the following form. For a
partition of unity \(\phi_i\) subordinate to charts and \(\psi_i\) as
before,
\[
 C(\tau)-\sum_i\phi_i\trOp_y\bigl(c^{(i)}(\cdot,\tau,\cdot)\bigr)\psi_i\in
 \cS\bigl(\R;\Psi^{-\infty}(H;E)\bigr),
\]
where each \(c^{(i)}(y,\tau,\eta)\) is a classical symbol of order \(\mu\) in
\(\zeta=(\tau,\eta)\in\R^{1+d}\). Holomorphic families are defined as in
Definition~\ref{tr:def:holM}.
\end{definition}

By \eqref{eq:Inkernel} these are the indicial families of
translation-invariant operators on \(\R_t\times H\) with rapidly decreasing
convolution kernels. Composition is pointwise in \(\tau\), and
\(\In([t,C])=-i\partial_\tau C\). We quote from (B1) and
\cite{MelroseEta}: \(\Psis^{\Z}\) is an algebra; \(\partial_\tau\) maps
\(\Psis^\mu\) to \(\Psis^{\mu-1}\); and if \(A\in\Psis^m\) is elliptic in
\((\tau,\eta)\) and \(A(\tau)\) is invertible for all \(\tau\), then
\(A^{-1}\in\Psis^{-m}\). For fixed \(\tau\), \(C(\tau)\in\Psi^\mu(H)\).

\begin{lemma}\label{tr:lem:susTr}
Let \(C(s)\in\Psis^{\mu(s)}(H;E)\) be holomorphic on \(\Omega\) with
\(\operatorname{Re}\mu<-d\). Then \(f(\tau;s)=\Tr C(s)(\tau)\) is a
holomorphic family of classical symbols on \(\R\) of order \(\mu(s)+d\):
\begin{gather*}
 f(\tau;s)\sim\sum_{j\ge0}f^\pm_j(s)|\tau|^{\mu(s)+d-j},\quad\tau\to\pm\infty,\\
 f^\pm_j(s)=(2\pi)^{-d}\sum_i\int\!\!\int\phi_i\tr
 c^{(i)}_{\mu(s)-j}(s;y,\pm1,\omega)\,d\omega\,dy .
\end{gather*}
The remainders are \(O(|\tau|^{\operatorname{Re}\mu+d-J})\), locally
uniformly in \(s\).
\end{lemma}

\begin{proof}
By Lemma~\ref{tr:lem:tc}, \(f=(2\pi)^{-d}\sum_i\int\!\!\int\phi_i\tr
c^{(i)}(s;y,\tau,\eta)\,d\eta\,dy\) plus a Schwartz function of \(\tau\).
Split \(c=\sum_{j<J}(1-\chi(\zeta))c_{\mu-j}+r_J\), where \(\chi\) is supported
in \(|\zeta|<\frac12\). For \(|\tau|\ge1\) substitute \(\eta=|\tau|\omega\).
By homogeneity, \(\int c_{\mu-j}(y,\tau,\eta)\,d\eta=|\tau|^{\mu-j+d}\int
c_{\mu-j}(y,\pm1,\omega)\,d\omega\). This integral converges since
\(\operatorname{Re}(\mu-j)<-d\). Finally
\(\int|r_J|\,d\eta\le C\int\trang{(\tau,\eta)}^{\operatorname{Re}\mu-J}
d\eta\le C'\trang\tau^{\operatorname{Re}\mu-J+d}\).
\end{proof}

\paragraph{Regularized integrals.}
Let \(\trP\) be the space of \(u\in\CI(\R)\) such that, for every \(N\),
\(u(\tau)=\sum c^\pm_{\gamma,l}|\tau|^\gamma(\log|\tau|)^l+O(|\tau|^{-N})\)
for \(\pm\tau\ge1\). The sum is finite, over \(\operatorname{Re}\gamma>-N\)
and \(l\in\N_0\). Using \(\int_1^T\tau^\gamma(\log\tau)^l\,d\tau\), one sees
that for \(u\in\trP\),
\[
 \int_{-T}^Tu\,d\tau=\sum d_{\gamma,l}T^\gamma(\log T)^l+c+o(1),\qquad
 T\to\infty,
\]
with a finite sum over \(\operatorname{Re}\gamma\ge0\), \((\gamma,l)\ne(0,0)\).
We put \(\trbint u\,d\tau:=c\).

\begin{lemma}\label{tr:lem:bint}
\begin{enumerate}
\item \(\trbint\) is linear. It equals \(\int\) on integrable elements of
\(\trP\). It vanishes on polynomials and on odd functions.
\item \(\trP\) is closed under multiplication by \(\tau\) and under
\(I:u\mapsto\int_0^\tau u\). Classical symbols of any order lie in
\(\trP\). The negative powers in the expansion of \(Iu\) are obtained by
integrating the expansion of \(u\) term by term.
\item Let \(u\in\CI(\R)\) with \(u^{(k)}\in\trP\) for \(0\le k\le p\). Let
\(e^\pm_k\) be the coefficient of \(|\tau|^{-1-k}\) (without logarithm) in
\(u^{(k)}\) at \(\pm\infty\). Then
\[
 \trbint u\,d\tau=\frac{(-1)^p}{p!}\trbint\tau^pu^{(p)}\,d\tau+
 \sum_{k=0}^{p-1}\frac{(-1)^k}{(k+1)!}\bigl(e^+_k+(-1)^ke^-_k\bigr).
\]
\end{enumerate}
\end{lemma}

\begin{proof}
(1) and (2) are immediate from the definitions. For the last part of (2),
\(\int_0^\tau=\int_0^1+\int_1^\tau\), and a constant of integration affects
only the power \(\tau^0\). For (3), integrate by parts \(p\) times on
\([-T,T]\):
\[
 \int_{-T}^Tu=\sum_{k<p}\frac{(-1)^k}{(k+1)!}\bigl[\tau^{k+1}u^{(k)}\bigr]_{-T}^T
 +\frac{(-1)^p}{p!}\int_{-T}^T\tau^pu^{(p)} .
\]
Now take constant terms. The constant term of \(T^{k+1}u^{(k)}(T)\) is
\(e^+_k\). That of \((-T)^{k+1}u^{(k)}(-T)\) is \((-1)^{k+1}e^-_k\).
\end{proof}

\begin{lemma}[Finite parts of holomorphic families]\label{tr:lem:FP}
Let \(f(\cdot;s)\) be a holomorphic family of classical symbols on \(\R\) of
order \(\mu_0-s\), \(\mu_0\in\Z\), on a connected open set \(\Omega\ni0\)
meeting \(\{\operatorname{Re}s>\mu_0+1\}\). Then \(\int f(\tau;s)\,d\tau\)
extends meromorphically to \(\Omega\), with at most simple poles at
\(s\in\mu_0+1-\N_0\). Let \(j_0=\mu_0+1\) and
\(e(s)=f^+_{j_0}(s)+f^-_{j_0}(s)\), with \(e=0\) if \(j_0<0\). Then
\[
 \Res_{s=0}\int f\,d\tau=e(0),\qquad
 \FP_{s=0}\int f\,d\tau=\trbint f(\tau;0)\,d\tau+e'(0).
\]
\end{lemma}

\begin{proof}
Write \(f=\sum_{j<J}\chi_\pm f^\pm_j(s)|\tau|^{\mu_0-s-j}+r_J\), where
\(\chi_\pm\) is the indicator of \(\{\pm\tau\ge1\}\) and \(r_J\) is
holomorphic with values in \(L^1\) for \(J>\operatorname{Re}\mu_0+2\). Since
\(\int_1^\infty\tau^{\mu_0-s-j}\,d\tau=(s-(\mu_0+1-j))^{-1}\), the
continuation is
\[
 \int_{-1}^1f+\sum_{j<J,\pm}\frac{f^\pm_j(s)}{s-(\mu_0+1-j)}+
 \int\Bigl(r_J-\chi_{[-1,1]}\sum_{j<J}(\ldots)\Bigr).
\]
Here the last integral is holomorphic. At \(s=0\) the term \(j_0\)
contributes \(e(s)/s=e(0)/s+e'(0)+O(s)\). The other terms are holomorphic
near \(0\). On the other hand, \(\trbint f(\cdot;0)\) is given by the same
expression at \(s=0\), with the \(j_0\) term replaced by \(0\), because
\(\int_1^T\tau^{-1}d\tau=\log T\) has constant term \(0\). For \(j\ne j_0\),
\(\int_1^T\tau^{\mu_0-j}\) has constant term \(-(\mu_0-j+1)^{-1}\), which is
the value at \(s=0\) of the continued term.
\end{proof}

\begin{definition}[Melrose's regularized trace]\label{tr:def:bTr}
For \(C\in\Psis^m(H;E)\), \(m\in\Z\), choose \(p\in\N_0\) with \(p\ge
m+d+1\) and put \(h_p(\tau)=\Tr\bigl(\partial_\tau^pC(\tau)\bigr)\). This
is a classical symbol of order \(m+d-p<0\) by Lemma~\ref{tr:lem:susTr}.
Define
\[
 \bTr(C)=\frac1{2\pi}\trbint I^ph_p\,d\tau .
\]
This is the definition \cite[(58)--(59)]{MN}: \(I^ph_p\) is the \(p\)-fold
integral from \(0\), and \(\trbint\) takes the constant term of
\(\int_{-T}^T\).
\end{definition}

\begin{theorem}[The trace property, (B7)]\label{tr:thm:bTr}
\begin{enumerate}
\item \(\bTr(C)\) does not depend on \(p\). If \(m<-d-1\), then
\(\bTr(C)=(2\pi)^{-1}\int\Tr C(\tau)\,d\tau\). In particular \(\bTr\) is the
functional of Lemma~\ref{lem:Trd} on \(\Psis^{-\infty}\).
\item If \(C(-\tau)=-C(\tau)\), then \(\bTr(C)=0\).
\item \(\bTr([A,B])=0\) for all \(A,B\in\Psis^{\Z}(H;E)\). Thus \(\bTr\) is
a trace on \(\Psis^{\Z}(H;E)\).
\end{enumerate}
\end{theorem}

\begin{proof}
(1) We have \(h_{p+1}=h_p'\), so \(I^{p+1}h_{p+1}=I^p(h_p-h_p(0))=
I^ph_p-h_p(0)\tau^p/p!\). A polynomial has \(\trbint=0\). If \(m<-d-1\),
then \(\Tr C(\tau)\) is integrable and differs from \(I^ph_p\) by its
Taylor polynomial of degree \(p-1\) at \(0\).

(2) Then \(h_p\) has parity \((-1)^{p+1}\). Each integration from \(0\)
reverses the parity, so \(I^ph_p\) is odd, and \(\int_{-T}^T\) of it
vanishes.

(3) Let \(A\in\Psis^{m_1}\) and \(B\in\Psis^{m_2}\), and take \(p>m_1+m_2+d\).
Since composition is pointwise in \(\tau\),
\[
 \partial_\tau^p[A,B](\tau)=\sum_{l=0}^p\binom pl\bigl[\partial_\tau^lA(\tau),
 \partial_\tau^{p-l}B(\tau)\bigr].
\]
Each commutator is a commutator of operators on \(H\) of orders adding to
\(m_1+m_2-p<-d\). By Lemma~\ref{tr:lem:tc}(2) its trace vanishes. So
\(h_p\equiv0\) and \(\bTr([A,B])=0\).
\end{proof}

The \emph{suspended residue} of \(C\in\Psis^m(H;E)\) is the Wodzicki
residue of \(C\), regarded as a translation-invariant operator on
\(\R\times H\), per unit length in \(t\):
\[
 \trTrt(C)=(2\pi)^{-(d+1)}\sum_i\int\phi_i\int_{S^d}\tr c^{(i)}_{-1-d}
 (y,\zeta)\,dS(\zeta)\,dy .
\]

\begin{theorem}[Holomorphic regularization of \(\bTr\)]\label{tr:thm:susreg}
Let \(C(s)\in\Psis^{m-s}(H;E)\), \(m\in\Z\), be holomorphic for
\(\operatorname{Re}s>-\frac12\). For instance \(C(s)=CQ_0(s)\) with
\(Q_0\) a holomorphic suspended family of order \(-s\) and
\(Q_0(0)=\mathrm{Id}\). Put
\[
 c^\pm(s)=(2\pi)^{-d}\sum_i\int\!\!\int\phi_i\tr c^{(i)}_{-1-d-s}(s;y,\pm1,
 \omega)\,d\omega\,dy,\qquad\gamma(s)=c^+(s)+c^-(s),
\]
where \(c_{-1-d-s}(s)\) is the homogeneous term of degree \(-1-d-s\) of the
symbol of \(C(s)\). Then \(G(s)=(2\pi)^{-1}\int\Tr C(s)(\tau)\,d\tau\),
defined for \(\operatorname{Re}s>m+d+1\), extends meromorphically to
\(\operatorname{Re}s>-\frac12\). Moreover
\[
 \Res_{s=0}G=\frac{\gamma(0)}{2\pi}=\trTrt(C(0)),\qquad
 \FP_{s=0}G=\bTr(C(0))+\frac{\gamma'(0)}{2\pi}.
\]
The correction \(\gamma'(0)\) is a local functional of the symbols of
\(C(s)\) near \(s=0\).
\end{theorem}

\begin{proof}
Fix \(p\ge m+d+2\). On \(\operatorname{Re}s>-\frac12\) the family
\(\partial_\tau^pC(s)\) has order with real part \(<-d\). By
Lemma~\ref{tr:lem:susTr}, \(h_p(\cdot;s)=\Tr\partial_\tau^pC(s)\) is a
holomorphic family of classical symbols of order \(m-s+d-p\). Its
homogeneous term of degree \(-1-d-s-p\) is \(\partial_\tau^pc_{-1-d-s}\).
Differentiation under the integral sign is legitimate since
\(\operatorname{Re}s>-1\). Hence the coefficient \(a^\pm(s)\) of
\(|\tau|^{-1-p-s}\) in \(h_p\) is the coefficient produced by
\(\partial_\tau^p(c^\pm(s)|\tau|^{-1-s})\):
\[
 a^+(s)=(-1)^p(1+s)_pc^+(s),\qquad a^-(s)=(1+s)_pc^-(s),\qquad
 (1+s)_p=\prod_{k=1}^p(k+s).
\]

For \(\operatorname{Re}s>m+d+1\) and \(s\) outside a discrete set, the
function \(f_s=\Tr C(s)\) is integrable and \(f_s-I^ph_p(\cdot;s)\) is a
polynomial. By Lemma~\ref{tr:lem:bint}(3) (with no boundary constants,
since no exponent equals \(-1-k\)),
\[
 G(s)=\frac1{2\pi}\frac{(-1)^p}{p!}\trbint\tau^ph_p(\tau;s)\,d\tau .
\]
The family \(\tau^ph_p(\cdot;s)\) is a holomorphic family of classical
symbols of order \(m+d-s\). The coefficient of \(|\tau|^{-1-s}\) in it is
\(\epsilon(s)=a^+(s)+(-1)^pa^-(s)=(-1)^p(1+s)_p\gamma(s)\). By
Lemma~\ref{tr:lem:FP},
\begin{gather*}
 \Res_{s=0}G=\frac{(1+s)_p\gamma(s)}{2\pi\,p!}\Big|_{s=0}=\frac{\gamma(0)}{2\pi},\\
 \FP_{s=0}G=\frac1{2\pi}\Bigl(\frac{(-1)^p}{p!}\trbint\tau^ph_p(\tau;0)\,d\tau
 +H_p\gamma(0)+\gamma'(0)\Bigr),
\end{gather*}
where \(H_p=\sum_{k=1}^p\frac1k\) is the derivative of \((1+s)_p/p!\) at
\(s=0\). Now evaluate \(\bTr(C(0))=(2\pi)^{-1}\trbint I^ph_p(\cdot;0)\) with
Lemma~\ref{tr:lem:bint}(3), taking \(u=I^ph_p\). The coefficient of
\(|\tau|^{-1}\) in \(u\) is \(b^+=(-1)^pa^+(0)/p!\) at \(+\infty\) and
\(b^-=a^-(0)/p!\) at \(-\infty\). Then \(e^+_k=(-1)^kk!\,b^+\) and
\(e^-_k=k!\,b^-\). The boundary sum is
\[
 H_p(b^++b^-)=H_p\tfrac{(-1)^p}{p!}\epsilon(0)=H_p\gamma(0).
\]
So
\(\FP_{s=0}G-\bTr(C(0))=\gamma'(0)/2\pi\).

Finally, let \(g\) be homogeneous of degree \(-(d+1)\) on
\(\R^{d+1}\setminus0\). The substitution \(\zeta=(\tau,\tau\omega)\) in
\(\int_{1\le|\zeta|\le e}g\,d\zeta=\int_{S^d}g\,dS\) gives
\(\int_{S^d}g\,dS=\int g(1,\omega)\,d\omega+\int g(-1,\omega)\,d\omega\).
Hence \(\gamma(0)=2\pi\,\trTrt(C(0))\).
\end{proof}

\begin{example}[\(\FP\ne\bTr\)]\label{tr:ex:counter}
Let \(H\) be a point, \(C=\trang\tau\) (the indicial family of
\(\trang{D_t}=Q_0\)), and \(Q_0(s)=\trang\tau^{-s}\). Then
\(\int_\R\trang\tau^{1-s}\,d\tau=\sqrt\pi\,\Gamma(\frac{s-2}2)/
\Gamma(\frac{s-1}2)=\frac1s+\log2-\frac12+O(s)\). On the other hand,
\(\int_{-T}^T\trang\tau\,d\tau=T^2+\log T+\log2+\frac12+o(1)\). So
\[
 \FP_{s=0}\,\frac1{2\pi}\int\trang\tau^{1-s}d\tau=\frac{\log2-\frac12}{2\pi}
 \ne\frac{\log2+\frac12}{2\pi}=\bTr(\trang\tau).
\]
This agrees with Theorem~\ref{tr:thm:susreg}. The coefficient of
\(|\tau|^{-1-s}\) in \(\trang\tau^{1-s}\) is \(\frac{1-s}2\) at each end,
so \(\gamma(s)=1-s\), \(\gamma(0)/2\pi=1/2\pi\) is the residue, and the
correction is \(\gamma'(0)/2\pi=-1/2\pi\).
\end{example}

\begin{correction}[MN Proposition 8, Corollary 2 and Proposition 17]
\label{tr:corr:Prop8}
The residue statement of \cite[Prop.~8]{MN}
(Proposition~\ref{prop:Tr8}) is correct. The claim that the regular value
of \(\bTr(BQ_0^{-s})\) at \(s=0\) is \(\bTr(B)\) is false in general
(Example~\ref{tr:ex:counter}). The correct statement is
Theorem~\ref{tr:thm:susreg} with \(C(s)=BQ_0(s)\). Correspondingly,
Corollary~\ref{cor:Tr61}(2) reads
\[
 \hTrd(A)=\bTr\bigl(A_{(1)}\bigr)+\frac{\gamma_{A_{(1)}}'(0)}{2\pi},
\]
where \(\gamma_{A_{(1)}}\) is computed from \(A_{(1)}Q_0(s)\). Here \(Q\)
is translation invariant to all orders at \(\partial X\), with Taylor series
\(Q_0(s)\). Indeed, by Lemma~\ref{lem:LR} and
Theorem~\ref{tr:thm:mero}(2), \(\hTrd(A)\) is the regular value at \(s=0\)
of \((2\pi)^{-1}\int\Tr\bigl(A_{(1)}Q_0(s)\bigr)\,d\tau\), because
\((AQ(s))_{(1)}=A_{(1)}Q_0(s)\). Proposition~\ref{prop:eta} then holds with
a local correction; see Proposition~\ref{tr:prop:etabar}.
\end{correction}

\begin{remark}[Parity (sketch)]\label{tr:rem:parity}
Say that a symbol of order \(\mu\) has \emph{parity type} \(\varepsilon\) if
\(c_{\mu-k}(y,-\zeta)=\varepsilon(-1)^kc_{\mu-k}(y,\zeta)\) for all \(k\).
Differential operators of order \(m\), their parametrices, \(\partial_\tau\)
of them, and the family \(\trang\zeta^{-s}\) have this property. The types
are \((-1)^m\), \((-1)^m\), \((-1)^{m-1}\) and \(+1\). The property is
preserved by composition and by changes of coordinates on \(H\), since the
terms of the composition and coordinate-change formulas are
\(\partial_\zeta^\alpha\) of symbols multiplied by polynomials in \(\zeta\)
of matching parity. For \(C=A_0^{-1}[t,A_0]\) with \(A_0\) the indicial
family of a differential operator, \(C\,Q_0(s)\) has type \(-1\). Its term
of degree \(-1-d-s\) has \(k=d\), so
\(c^+(s)+c^-(s)=(1+(-1)^{d+1})\int c(1,\omega)\,d\omega\). This vanishes
when \(d\) is even, that is, when \(n=\dim X\) is \emph{odd}. In that case
\(\gamma\equiv0\). In the APS case, \(n\) even, parity alone does not kill
\(\gamma\).
\end{remark}

\subsection{The eta invariant and Dirac operators}\label{tr:subsec:eta}

\begin{definition}\label{tr:def:eta}
Let \(A\in\Psis^m(H;E,F)\) be elliptic and invertible. Then
\(A^{-1}\partial_\tau A\in\Psis^{-1}(H;E)\), and Melrose's eta invariant
\cite{MelroseEta} is
\[
 \eta(A)=\frac1{\pi i}\overline{\int_\R}\Tr\bigl(A^{-1}\partial_\tau A
 \bigr)\,d\tau:=\frac{2\pi}{\pi i}\,\bTr\bigl(A^{-1}\partial_\tau A\bigr)
 =2\,\bTr\bigl(A^{-1}[t,A]\bigr).
\]
The last equality uses \(\In([t,A])=-i\partial_\tau A\) from
\eqref{eq:Inkernel}. The notation \(\overline{\int}\Tr(\cdot)\,d\tau\) always
means \(2\pi\bTr\). The pointwise trace exists only for \(d=0\).
\end{definition}

\begin{theorem}[Dirac operators, (B7)]\label{tr:thm:Dirac}
Let \(B\) be a self-adjoint, invertible, first-order elliptic differential
operator on \(H\). Let \(\sigma:E\to F\) be a unitary bundle isomorphism,
and let \(A_0(\tau)=\sigma(i\tau+B)\). This is the indicial family of
\(\sigma(-\partial_t+B)\), the APS normal form \(\sigma(\partial_u+B)\) in
\(u=-t\) of Corollary~\ref{cor:Dirac}. Then
\[
 \eta(A_0)=\eta_B(0),\qquad\eta_B(s)=\sum_{\lambda\in\operatorname{spec}B}
 \sgn\lambda\,|\lambda|^{-s}.
\]
\end{theorem}

\begin{proof}
We have \([t,A_0]=-i\partial_\tau A_0=\sigma\), so
\(A_0^{-1}[t,A_0]=(B+i\tau)^{-1}=:C\in\Psis^{-1}(H;E)\). Split
\(C=C_e+C_o\) with \(C_e(\tau)=B(B^2+\tau^2)^{-1}\) and
\(C_o(\tau)=-i\tau(B^2+\tau^2)^{-1}\). Both lie in \(\Psis^{-1}\). By
Theorem~\ref{tr:thm:bTr}(2), \(\bTr(C)=\bTr(C_e)\). We must show
\(\bTr(C_e)=\frac12\eta_B(0)\).

Let \(\theta(t)=\Tr(Be^{-tB^2})=\sum_\lambda\lambda e^{-t\lambda^2}\). It
decays exponentially as \(t\to\infty\). From \cite[Section ``The eta invariant'']{Notes}
we quote the small-time expansion \(\theta(t)\sim\sum_{k\ge0}a_kt^{(k-d-1)/2}\)
as \(t\downarrow0\). We also quote the Mellin representation
\(\Gamma(\frac{s+1}2)\eta_B(s)=\int_0^\infty t^{(s-1)/2}\theta(t)\,dt\)
and the regularity of \(\eta_B\) at \(0\), which is equivalent to \(a_d=0\).
Split \(\theta=\sum_{k<K}a_kt^{(k-d-1)/2}+r_K\) on \((0,1]\), with
\(K>d+1\). Then
\begin{equation}\label{tr:eq:etaMellin}
 \sqrt\pi\,\eta_B(0)=\int_1^\infty\frac{\theta}{\sqrt t}\,dt+\int_0^1
 \frac{r_K}{\sqrt t}\,dt+\sum_{k<K,\,k\ne d}\frac{2a_k}{k-d}.
\end{equation}

Fix \(q>(d+1)/2\), put \(p=2q\), and let
\(E_q(x)=e^{-x}-\sum_{k<q}(-x)^k/k!\). Define
\(F(\tau)=\int_0^\infty\theta(t)E_q(t\tau^2)\,dt\). The integral converges:
near \(t=0\), \(E_q=O(t^q)\) and \(\theta=O(t^{-(d+1)/2})\). The polynomial
part of \(E_q(t\tau^2)\) has degree \(<p\) in \(\tau\). Differentiation
under the integral and Fubini over the eigenvalues give
\[
 F^{(p)}(\tau)=\int_0^\infty\theta(t)\,\partial_\tau^pe^{-t\tau^2}\,dt=
 \sum_\lambda\partial_\tau^p\frac{\lambda}{\lambda^2+\tau^2}=\Tr\bigl(\partial
 _\tau^pC_e(\tau)\bigr)=h_p .
\]
The series converges absolutely because \(\sum|\lambda|^{-p-1}<\infty\). So
\(F-I^ph_p\) is a polynomial, and \(\bTr(C_e)=(2\pi)^{-1}\trbint F\,d\tau\).

By Fubini and the substitution \(u=\sqrt t\,\tau\),
\[
 \int_{-T}^TF\,d\tau=\int_0^\infty\theta(t)\,t^{-1/2}W(T\sqrt t)\,dt,
 \qquad W(R)=\int_{-R}^RE_q(u^2)\,du=\sqrt\pi-\sum_{l<q}c_lR^{2l+1}+e(R),
\]
with \(c_l=2(-1)^l/(l!(2l+1))\) and \(e(R)=-2\int_R^\infty e^{-u^2}du\).
The function \(e\) is bounded and decays faster than any power. We find the
constant term as \(T\to\infty\) piece by piece.
\begin{itemize}
\item On \([1,\infty)\):
\(\int_1^\infty\theta\,t^{-1/2}W(T\sqrt t)=\sqrt\pi\int_1^\infty\theta
t^{-1/2}-\sum_lc_lT^{2l+1}\int_1^\infty\theta t^l+o(1)\). The constant term
is \(\sqrt\pi\int_1^\infty\theta t^{-1/2}\).
\item The remainder \(r_K\) on \((0,1]\) gives
\(\sqrt\pi\int_0^1r_Kt^{-1/2}\) by the same argument. The \(e\)-term tends
to \(0\) by dominated convergence.
\item For \(a_kt^\beta\), \(\beta=(k-d-1)/2\), substitute \(R=T\sqrt t\):
\[
 \int_0^1t^{\beta-1/2}W(T\sqrt t)\,dt=2T^{-2\beta-1}\int_0^TR^{2\beta}W(R)\,dR .
\]
Insert \(W=\sqrt\pi-\sum_lc_lR^{2l+1}+e\) on \([1,T]\). If \(k\ne d\), the
only constant term is \(2\sqrt\pi/(2\beta+1)=2\sqrt\pi/(k-d)\). The other
terms are non-zero powers of \(T\), possibly times \(\log T\), or
\(o(1)\). The term \(k=d\) is absent because \(a_d=0\).
\end{itemize}
Comparing with \eqref{tr:eq:etaMellin}, \(\trbint F\,d\tau=\sqrt\pi\cdot
\sqrt\pi\,\eta_B(0)=\pi\eta_B(0)\). Hence \(\bTr(C)=\frac12\eta_B(0)\) and
\(\eta(A_0)=2\bTr(C)=\eta_B(0)\).
\end{proof}

The proof also shows the formal identity
\(\pi^{-1}\int(\lambda+i\tau)^{-1}d\tau=\sgn\lambda\) of
Corollary~\ref{cor:Dirac}, now with the regularization made precise. As a
numerical check, take \(H=S^1\) and \(B=D_\theta+a\), \(0<a<1\). The partial
integrals \(\int_{-T}^TF\) computed from the spectrum \(\Z+a\) have constant
term \(\pi(1-2a)=\pi\eta_B(0)\).

\begin{proposition}[The functional \(\etab\) of MN (89)]\label{tr:prop:etabar}
Let \(A\in\Psic^m(X;E,F)\) be fully elliptic, with \(A_0=\In(A)\), and let
\(Q\) be admissible and translation invariant to all orders at
\(\partial X\), with Taylor series \(Q_0(s)\) (on \(E\) and on \(F\)).
\begin{enumerate}\renewcommand{\labelenumi}{(\roman{enumi})}
\item \(\etab(A)=\eta(A_0)+\delta_Q(A)\), where
\(\delta_Q(A)=(2\pi)^{-1}\gamma'(0)\) and \(\gamma\) is computed in
Theorem~\ref{tr:thm:susreg} from \(A_0^{-1}[t,A_0]Q_0(s)\) and
\([t,A_0]A_0^{-1}Q_0(s)\). The correction \(\delta_Q(A)\) is a local
functional of the boundary symbols. By Remark~\ref{tr:rem:parity} it
vanishes when \(n\) is odd and \(A\) is a differential operator (sketch).
\item In the situation of Theorem~\ref{tr:thm:Dirac}, suppose that
\(Q_0(s)=(D_t^2+B^2)^{-s/2}\) on \(E\) and
\(\sigma(D_t^2+B^2)^{-s/2}\sigma^{-1}\) on \(F\). Then
\(\etab(A)=\eta(A_0)=\eta_B(0)\). Such \(Q\) exist if one quotes the fact
that \((D_t^2+B^2)^{-s/2}\) is an entire family in \(\Psis^{-s}(H)\)
(Seeley's complex powers \cite{Seeley} with the parameter \(\tau\),
quoted: \cite[App.~B]{MN}). Then take
\(Q(s)=\tilde\omega Q_0(s)\omega+Q_X(s)(1-\omega)\), with \(Q_X\) from
Proposition~\ref{tr:prop:Q}.
\end{enumerate}
\end{proposition}

\begin{proof}
(i) As in the proof of Proposition~\ref{prop:eta},
\((A^{-1}[\lx,A])_{(1)}=-A_0^{-1}[t,A_0]=:-C_1\) and
\(([\lx,A]A^{-1})_{(1)}=-[t,A_0]A_0^{-1}=:-C_2\). By
Correction~\ref{tr:corr:Prop8}, \(\etab(A)=-\hTrd(A^{-1}[\lx,A]+[\lx,A]
A^{-1})=\bTr(C_1)+\bTr(C_2)+(2\pi)^{-1}\gamma'_{C_1+C_2}(0)\). By
Theorem~\ref{tr:thm:bTr}(3), \(\bTr(C_2)=\bTr(C_1)\), and
\(2\bTr(C_1)=\eta(A_0)\).

(ii) We have \(C_1=(B+i\tau)^{-1}\) and \(C_2=\sigma C_1\sigma^{-1}\), so
\(\Tr(C_2Q_0^F(s))=\Tr(C_1Q_0^E(s))\). For \(\operatorname{Re}s>d\), the
eigenvalue sum and the \(\tau\)-integral converge absolutely. The odd part
integrates to \(0\), and
\[
 \frac1{2\pi}\int\Tr\bigl(C_1(B^2+\tau^2)^{-s/2}\bigr)d\tau=\frac1{2\pi}
 \sum_\lambda\lambda\int_\R(\lambda^2+\tau^2)^{-1-s/2}d\tau=
 \frac{\Gamma(\frac{1+s}2)}{2\sqrt\pi\,\Gamma(1+\frac s2)}\,\eta_B(s).
\]
This is regular at \(s=0\) with value \(\frac12\eta_B(0)\). By
Correction~\ref{tr:corr:Prop8} (before the comparison with \(\bTr\)),
\(\hTrd(A^{-1}[\lx,A])=\hTrd([\lx,A]A^{-1})=-\frac12\eta_B(0)\). So
\(\etab(A)=\eta_B(0)\), which is \(\eta(A_0)\) by
Theorem~\ref{tr:thm:Dirac}. For the construction of \(Q\): the family is
admissible since \(\tilde\omega\omega=\omega\), and the term
\(Q_X(s)(1-\omega)\) has kernel supported away from the front face.
\end{proof}

\begin{remark}[The regularizer \((A^*A+1)^{1/2}\) (sketch)]\label{tr:rem:plusone}
Theorem~\ref{thm:reduced} uses \(Q=(A^*A+1)^{1/2m}\), whose indicial family
is \((D_t^2+B^2+1)^{1/2}\) for \(A\) as in Theorem~\ref{tr:thm:Dirac}.
Expanding \((1+(\tau^2+B^2)^{-1})^{-s/2}\) binomially in the computation
above (for \(|\lambda|>1\)) gives
\[
 \etab(A)=\eta_B(0)+\frac{1}{\sqrt\pi}\sum_{1\le k\le d/2}\frac{(-1)^k}
 {k\cdot k!}\,a_{d-2k},
\]
where the \(a_j\) are the heat coefficients of \(\theta(t)\). For Dirac
operators compatible with a Clifford structure,
\(\theta(t)=O(t^{1/2})\) as \(t\downarrow0\) (local cancellation of
Bismut--Freed \cite{BismutFreed}, quoted). So all \(a_j\) with \(j\le d\) vanish, and
\(\etab(A)=\eta_B(0)\) also for this \(Q\). For general first-order \(B\)
the extra terms need not vanish. They are local and are compensated in the
index formula by \(\beta_Q\) (Proposition~\ref{tr:prop:transfer}(3)).
\end{remark}

\begin{remark}[Signs]\label{tr:rem:signs}
In \cite{Notes} the APS operator is \(P=c(du)(\partial_u+A^+)\) near the
boundary, with \(u=\log x_b\to-\infty\) on the end, and
\(\operatorname{ind}P=\int\hat A\operatorname{ch}-\frac12(\eta(A^+)+h)\).
The cusp picture has \(x_b=e^{-1/x}\), so \(u=-t\) and
\(\partial_u=-\partial_t\). Hence \(A=\sigma(-\partial_t+B)\) with
\(\sigma=c(du)\) and \(B=A^+\), and Theorem~\ref{tr:thm:Dirac} gives
\(\eta(A_0)=\eta(A^+)\). Thus the corrected formula
\(\Ind A=\AS(A)-\frac12\eta(A_0)\) of Theorem~\ref{thm:reduced}, in which
\(\delta_Q=0\) here (Remark~\ref{tr:rem:plusone}), reproduces
APS (with \(h=0\)). The sign of the eta term is fixed by the indicial
convention \eqref{eq:Indef}, \(\In(A,\tau)=e^{i\tau t}Ae^{-i\tau t}\); with
\(e^{-i\tau t}Ae^{i\tau t}\) instead, \(\eta(A_0)\) would change sign. This
confirms the sign of the eta term in Correction~\ref{corr:ASsign}.
\end{remark}
